\documentclass[10pt,a4paper]{article}
\usepackage{amsmath}
\newcommand{\lab}{the laboratory}
\title{Results of the Seed Germination Experiment}
\author{Anselm Brodrick}
\date{April 18, 2019}
\begin{document}
\maketitle

\section{Purpose}
The aim of this experiment was to compare germination rates of three kinds of
seed under identical conditions in \lab. The trays were kept at a constant
temperature and watered every second day.

\section{Results}
Each tray held one hundred seeds. After two weeks the following counts were
recorded.

\begin{table}[htbp]
\centering
\begin{tabular}{|l|c|p{3cm}|}
\hline
Variety & Germinated & Remarks \\
\hline
Basil & 82 & Uniform emergence, strong seedlings \\
\hline
Parsley & 47 & Slow and uneven, some seeds mouldy \\
\hline
Coriander & 91 & Very rapid germination \\
\hline
\end{tabular}
\caption{Germination counts after fourteen days.}\label{tab:germ}
\end{table}

The mean rate is $\bar{x} = \frac{1}{3}\sum_{i=1}^{3} x_i = 73.3$ per hundred.
The sample variance is
\[
s^2 = \frac{1}{2}\sum_{i=1}^{3} (x_i - \bar{x})^2 \approx 485.
\]

\begin{figure}[htbp]
\centering
\begin{tabular}{lcr}
Day 3 & Day 7 & Day 14 \\
12 & 40 & 73 \\
\end{tabular}
\caption{Average number of germinated seeds per tray over time.}
\label{fig:time}
\end{figure}

\section{Discussion}
As Table~\ref{tab:germ} shows, parsley performed notably worse. Figure~\ref{fig:time}
on page~\pageref{fig:time} indicates that growth is roughly linear early on.

\newpage
\section{Conclusion}
Coriander and basil are recommended for the school garden. Further trials
with different watering schedules would clarify whether the poor parsley
result is intrinsic or a consequence of the chosen conditions.
\end{document}
